\chapter{Proposed Approach}
\label{chap:proposed-approach}

% PLAN: draft in W9--W10 (holiday buffer), ~5--8 pages. Provisional by design:
% the intermediate report must show a justified direction, not final results.

\section{Problem Formulation}
\label{sec:pa-formulation}

% Formal fair recourse problem using the notation of Section~\ref{sec:bg-notation}:
% classifier, cost, actions, fairness constraint/objective on recourse.

\section{Candidate Architectures}
\label{sec:pa-architectures}

% Decision matrix from Section~\ref{sec:lr-gaps}:
% O1 KAN as decision model | O2 KAN as causal model | O3 KAN as recourse
% generator | O4 combination. Criteria: feasibility, novelty, evaluability,
% interpretability. State the recommended option.

\section{Experimental Design}
\label{sec:pa-experimental-design}

% Datasets table (from the Evaluation tab: \cite{lequy2022survey, becker1996adult,
% hofmann1994statlog, angwin2016machine, wightman1998lsac, yeh2009comparisons};
% benchmark \cite{pawelczyk2021carla}; tool \cite{wexler2020whatif}), baselines (MLP counterparts + recourse
% methods from Sections~\ref{sec:lr-cf-search} and~\ref{sec:lr-plausibility}), metrics
% (predictive, recourse, recourse-fairness), statistical tests, protocol.
% Optional: preliminary sanity experiment, if agreed with the supervisor.

\section{Summary}
\label{sec:pa-summary}
